\documentclass[a4paper,12pt]{article}

% Pacotes básicos
\usepackage{iftex}
\ifPDFTeX
    \usepackage[utf8]{inputenc}
    \usepackage[T1]{fontenc}
    \usepackage{lmodern}
    \usepackage[portuguese]{babel}
\else
    \usepackage{fontspec}
    \setmainfont{Latin Modern Roman}
    \usepackage{polyglossia}
    \setmainlanguage{portuguese}
\fi

\usepackage{graphicx}
\usepackage{geometry}
\usepackage{amsmath}
\usepackage{hyperref}
\usepackage{float}
\usepackage{caption}
\usepackage{subcaption}
\usepackage{listings}
\usepackage{xcolor}
\usepackage{setspace}
\usepackage{tabularx}
\usepackage{booktabs}
\usepackage{array}
\usepackage{indentfirst}

% Recuo da primeira linha, inclusive após títulos de seções
\setlength{\parindent}{1.25cm}
% Flexibilidade adicional para acomodar termos técnicos sem ultrapassar margens
\setlength{\emergencystretch}{2em}

% Coluna de largura ajustável com quebras de linha e alinhamento à esquerda
\newcolumntype{Y}{>{\raggedright\arraybackslash}X}

% Configuração das margens
\geometry{top=3cm,bottom=2cm,left=3cm,right=2cm}

% Comando para adicionar a Fonte nas figuras (Padrão ABNT)
\newcommand{\fonte}[1]{
    \vspace{-0.2cm} \begin{center} \small \textbf{Fonte:} #1 \end{center}
}

% Configuração para códigos
\definecolor{codegreen}{rgb}{0,0.6,0}
\definecolor{codegray}{rgb}{0.5,0.5,0.5}
\definecolor{codepurple}{rgb}{0.58,0,0.82}
\definecolor{backcolour}{rgb}{0.95,0.95,0.92}

\lstdefinestyle{mystyle}{
    backgroundcolor=\color{backcolour}, commentstyle=\color{codegreen},
    keywordstyle=\color{magenta}, numberstyle=\tiny\color{codegray},
    stringstyle=\color{codepurple}, basicstyle=\ttfamily\footnotesize,
    breakatwhitespace=false, breaklines=true, captionpos=b, keepspaces=true,
    numbers=left, numbersep=5pt, showspaces=false, showstringspaces=false,
    showtabs=false, tabsize=2
}

\lstset{style=mystyle}
\title{\textbf{Especificação Algébrica de Sistema de Biblioteca em CafeOBJ}}
\author{Seu Nome Aqui}
\date{Natal, \today}

\begin{document}

\begin{center}
    \thispagestyle{empty}
    \includegraphics[scale=0.05]{Imagens/ufrn-brasao.png}

    \textbf{\large Universidade Federal do Rio Grande do Norte}\\
    \textbf{\large Estruturas Algbricas DIM0167 - 2026.2} \\
    [1cm]
    \vspace*{2cm}
    {\Huge {Especificação Algébrica de Sistema de Biblioteca em CafeOBJ}}
    \\[0.5cm]
    \large Amauri Lacerda de Brito Junior \\
    \large Kael Amir Pina Paraguassu Silva \\
    \large Gabriel Victor da Silva \\
    \large Pedro Lucas Medeiros do Nascimento \\
    \large Samuel Morais da Rocha e Silva \\
    [1cm]
    \vfill
    \normalsize Natal, \today.
\end{center}

\newpage

\section{Introdução}

Este trabalho apresenta uma especificação algébrica de um sistema de biblioteca
em CafeOBJ. O modelo descreve livros, usuários e empréstimos por meio de tipos,
operadores e equações, com foco nas regras de negócio e na consistência do estado.
A validação experimental utiliza redução de termos e testes de operações válidas
e inválidas.

\section{Visão Geral do Sistema e Regras de Negócio}

O sistema de biblioteca é composto por três coleções principais:
\begin{itemize}
    \item Acervo: multiconjunto de livros pertencentes à biblioteca.
    \item Usuários: multiconjunto de usuários registrados.
    \item Empréstimos: multiconjunto de pares (livro, usuário) que representam
          empréstimos ativos.
\end{itemize}

Cada livro possui um identificador único (LivroId), uma obra associada (ObraId),
título, autor e ano. Cada usuário possui identificador único (UsuarioId), nome e
e-mail. As regras de negócio definidas são:

\begin{enumerate}
    \item Um livro só pode ser emprestado se pertencer ao acervo.
    \item Um livro não pode ser emprestado se já estiver emprestado.
    \item A devolução de um livro deve torná-lo novamente disponível.
    \item Um usuário só pode devolver um livro que esteja emprestado a ele.
    \item A remoção de um livro só é permitida quando ele estiver disponível.
    \item As identificações de livros e usuários devem ser únicas em suas
          respectivas coleções.
    \item Em um estado válido, a quantidade de livros disponíveis mais a
          quantidade de livros emprestados deve ser igual ao total do acervo.
    \item Um empréstimo exige que o usuário esteja registrado.
    \item Um usuário só pode ser removido se não possuir empréstimos ativos.
    \item O cadastro exige livro com título e autor não vazios e ano positivo,
          ou usuário com nome e e-mail não vazios, além de ID ainda não cadastrado.
\end{enumerate}

As transições globais aplicam essas condições por meio de guardas; quando uma
condição falha, preservam o estado recebido. Os predicados \texttt{valido?} e
\texttt{emprestimosValidos?} permitem consultar a consistência do estado.
Um usuário pode possuir vários exemplares simultaneamente. \texttt{LivroId}
identifica o exemplar físico, enquanto \texttt{ObraId} pode ser compartilhado por
exemplares da mesma obra. O catálogo oferece busca por ID, título e autor, além
da contagem de exemplares por obra.

\section{Assinatura}

\subsection{\textit{Sorts} e \textit{Subsorts}}

A especificação declara os seguintes \textit{sorts} e \textit{subsorts}:
\begin{itemize}
    \item \texttt{LivroId}, \texttt{UsuarioId}, \texttt{ObraId}: identificadores
          atômicos.
    \item \texttt{Livro}, \texttt{Usuario}: entidades.
    \item \texttt{Acervo}, \texttt{ConjuntoUsuarios},
          \texttt{ConjuntoEmprestimos}: coleções monoidais.
    \item \texttt{Emprestimo}: relação livro–usuário.
    \item \texttt{Biblioteca}: estado global visível.
    \item \texttt{LivroOuErro}: \textit{supersort} que inclui Livro e os
          construtores de erro \texttt{livroInexistente} e
          \texttt{livroDuplicado}.
    \item Subsorts: Livro < LivroOuErro, Livro < Acervo, Usuario <
          ConjuntoUsuarios, Emprestimo < ConjuntoEmprestimos.
\end{itemize}

\subsection{Operadores e Construtores}

As tabelas abaixo agrupam a assinatura por módulo. C indica construtor
(\texttt{constr}); D indica operador definido por equações. A sequência à
esquerda de $\to$ informa os tipos dos argumentos e sua quantidade; uma
sequência vazia indica aridade zero. Predicados retornam \texttt{Bool}.
Para compactar as tabelas, utilizam-se as abreviações de tipos:
\(LI=\texttt{LivroId}\), \(UI=\texttt{UsuarioId}\),
\(OI=\texttt{ObraId}\), \(L=\texttt{Livro}\),
\(U=\texttt{Usuario}\), \(A=\texttt{Acervo}\),
\(US=\texttt{ConjuntoUsuarios}\), \(ES=\texttt{ConjuntoEmprestimos}\),
\(B=\texttt{Biblioteca}\), \(R=\texttt{LivroOuErro}\),
\(S=\texttt{String}\), \(N=\texttt{Nat}\) e \(Q=\texttt{Bool}\).
Os tipos \texttt{String}, \texttt{Nat} e \texttt{Bool} são importados da
biblioteca padrão. \texttt{IDENTIFICADORES} declara apenas os três tipos de IDs;
as constantes concretas são introduzidas nos exemplos.

\begin{center}
\small
\begin{tabularx}{\linewidth}{@{}Y l c@{}}
\toprule
Operador (módulos LIVRO e USUARIO) & Assinatura & Tipo \\
\midrule
\texttt{livro} & $LI\;OI\;S\;S\;N\to L$ & C \\
\texttt{usuario} & $UI\;S\;S\to U$ & C \\
\texttt{id} & $L\to LI$; $U\to UI$ & D \\
\texttt{obra}; \texttt{ano} & $L\to OI$; $L\to N$ & D \\
\texttt{titulo}, \texttt{autor} & $L\to S$ & D \\
\texttt{nome}, \texttt{email} & $U\to S$ & D \\
\texttt{valido?} & $L\to Q$; $U\to Q$ & D \\
\bottomrule
\end{tabularx}
\end{center}

\begin{center}
\small
\begin{tabularx}{\linewidth}{@{}Y l c@{}}
\toprule
Operador (módulo ACERVO) & Assinatura & Tipo \\
\midrule
\texttt{vazioAcervo}; \texttt{\_\_} & $\to A$; $A\;A\to A$ & C \\
\texttt{livroInexistente}, \texttt{livroDuplicado} & $\to R$ & C \\
\texttt{iniciarAcervo} & $\to A$ & D \\
\texttt{adicionarLivro} & $L\;A\to A$ & D \\
\texttt{removerDoAcervo}, \texttt{filtrarId} & $LI\;A\to A$ & D \\
\texttt{buscarLivro}; \texttt{resultadoBusca} & $LI\;A\to R$; $A\to R$ & D \\
\texttt{buscarTitulo}, \texttt{tituloDeAutor} & $S\;A\to A$ & D \\
\texttt{qtdAcervo}; \texttt{qtdExemplaresObra} & $A\to N$; $OI\;A\to N$ & D \\
\texttt{pertenceAcervo?} & $LI\;A\to Q$ & D \\
\texttt{idsUnicos?}, \texttt{todosValidos?} & $A\to Q$ & D \\
\bottomrule
\end{tabularx}
\end{center}

\begin{center}
\small
\begin{tabularx}{\linewidth}{@{}Y l c@{}}
\toprule
Operador (módulo COLECAO-USUARIOS) & Assinatura & Tipo \\
\midrule
\texttt{vazioUsuarios}; \texttt{\_\_} & $\to US$; $US\;US\to US$ & C \\
\texttt{iniciarColecao} & $\to US$ & D \\
\texttt{adicionarUsuario} & $U\;US\to US$ & D \\
\texttt{removerUsuario} & $UI\;US\to US$ & D \\
\texttt{qtdUsuarios} & $US\to N$ & D \\
\texttt{pertence?} & $UI\;US\to Q$ & D \\
\texttt{idsUnicos?}, \texttt{todosValidos?} & $US\to Q$ & D \\
\midrule
Operador (módulo EMPRESTIMO) & Assinatura & Tipo \\
\midrule
\texttt{semEmprestimo}; \texttt{\_\_} & $\to ES$; $ES\;ES\to ES$ & C \\
\texttt{emprestar} (registro de empréstimo) & $LI\;UI\to ES$ & C \\
\texttt{iniciarEmprestimos} & $\to ES$ & D \\
\texttt{emprestar} (inserção na coleção) & $LI\;UI\;ES\to ES$ & D \\
\texttt{devolver} & $LI\;ES\to ES$ & D \\
\texttt{qtdEmprestimos} & $ES\to N$ & D \\
\texttt{emprestado?} & $LI\;ES\to Q$ & D \\
\texttt{temLivro?} & $UI\;LI\;ES\to Q$; $UI\;ES\to Q$ & D \\
\texttt{livrosUnicos?} & $ES\to Q$ & D \\
\bottomrule
\end{tabularx}
\end{center}

\begin{center}
\small
\begin{tabularx}{\linewidth}{@{}Y l c@{}}
\toprule
Operador (módulo BIBLIOTECA) & Assinatura & Tipo \\
\midrule
\texttt{biblioteca}; \texttt{bibliotecaVazia} & $A\;US\;ES\to B$; $\to B$ & C; D \\
\texttt{acervo}; \texttt{usuarios}; \texttt{emprestimos} & $B\to A$; $B\to US$; $B\to ES$ & D \\
\texttt{registrarLivro}; \texttt{removerLivro} & $L\;B\to B$; $LI\;B\to B$ & D \\
\texttt{registrarUsuario}; \texttt{removerUsuario} & $U\;B\to B$; $UI\;B\to B$ & D \\
\texttt{emprestarLivro}, \texttt{devolverLivro} & $LI\;UI\;B\to B$ & D \\
\texttt{qtdAcervo}, \texttt{qtdUsuarios}, \texttt{qtdEmprestimos}, \texttt{qtdDisponiveis} & $B\to N$ & D \\
\texttt{buscarLivro} & $LI\;B\to R$ & D \\
\texttt{buscarTitulo}, \texttt{tituloDeAutor} & $S\;B\to A$ & D \\
\texttt{qtdExemplaresObra} & $OI\;B\to N$ & D \\
\texttt{contarDisponiveis} & $A\;ES\to N$ & D \\
\texttt{pertenceAcervo?}, \texttt{emprestado?}, \texttt{disponivel?} & $LI\;B\to Q$ & D \\
\texttt{usuarioRegistrado?} & $UI\;B\to Q$ & D \\
\texttt{temLivro?} & $UI\;LI\;B\to Q$; $UI\;B\to Q$ & D \\
\texttt{emprestimosValidos?}, \texttt{valido?} & $B\to Q$ & D \\
\bottomrule
\end{tabularx}
\end{center}

\subsection{Variáveis}

As equações utilizam variáveis tipadas, quantificadas universalmente em seus
respectivos módulos. A tabela reúne as variáveis empregadas neste relatório;
alguns nomes locais foram uniformizados para evitar ambiguidades entre módulos.
As abreviações de tipos das tabelas anteriores não constituem declarações de
variáveis em CafeOBJ.

\begin{center}
\small
\begin{tabular}{@{}ll@{}}
\toprule
Declaração CafeOBJ & Papel \\
\midrule
\texttt{vars L L' : Livro .} & Exemplares \\
\texttt{var U : Usuario .} & Usuário \\
\texttt{vars LID LID2 : LivroId .} & Identificadores de exemplares \\
\texttt{var UID : UsuarioId .} & Identificador de usuário \\
\texttt{var OID : ObraId .} & Identificador de obra \\
\texttt{vars T AUT NOM EMAIL : String .} & Atributos textuais \\
\texttt{var ANO : Nat .} & Ano de publicação \\
\texttt{var A : Acervo .} & Coleção de livros \\
\texttt{var US : ConjuntoUsuarios .} & Coleção de usuários \\
\texttt{var ES : ConjuntoEmprestimos .} & Coleção de empréstimos \\
\texttt{var B : Biblioteca .} & Estado global \\
\bottomrule
\end{tabular}
\end{center}

\section{Estruturas de Dados}

Os módulos de entidades utilizam \texttt{IDENTIFICADORES}; \texttt{ACERVO}
utiliza \texttt{LIVRO}, \texttt{COLECAO-USUARIOS} utiliza \texttt{USUARIO}, e
\texttt{EMPRESTIMO} relaciona IDs de exemplares e usuários. O módulo
\texttt{BIBLIOTECA} integra as três coleções em um estado visível.

Cada coleção é um monóide comutativo. O operador binário \texttt{\_\_}
representa composição por justaposição: \texttt{L A}, por exemplo, reúne um
livro e um acervo. Seus atributos \texttt{assoc}, \texttt{comm} e \texttt{id:}
expressam associatividade, comutatividade e elemento neutro. Os neutros são
\texttt{vazioAcervo}, \texttt{vazioUsuarios} e \texttt{semEmprestimo}.
Assim, agrupamento e ordem não alteram a coleção; repetições são conservadas
e contadas. As inclusões por subsorts permitem usar um \texttt{Livro} ou um
\texttt{Usuario} como coleção unitária. O construtor de empréstimo produz
diretamente um \texttt{ConjuntoEmprestimos}.

O estado global é a tripla \texttt{biblioteca(\allowbreak{}A,\allowbreak{} US,\allowbreak{} ES)}. Suas projeções
recuperam os componentes sem modificá-los:
\begin{lstlisting}
eq bibliotecaVazia = biblioteca(vazioAcervo, vazioUsuarios, semEmprestimo) .
eq acervo(biblioteca(A, US, ES)) = A .
eq usuarios(biblioteca(A, US, ES)) = US .
eq emprestimos(biblioteca(A, US, ES)) = ES .
\end{lstlisting}

\section{Equações}

Em CafeOBJ, \texttt{eq} declara uma igualdade usada na redução de termos;
\texttt{if ... then ... else ... fi} seleciona o resultado conforme uma
condição booleana. As projeções de entidades recuperam os campos de seus
construtores, como \texttt{id(\allowbreak{}livro(\allowbreak{}LID,\allowbreak{} OID,\allowbreak{} T,\allowbreak{} AUT,\allowbreak{} ANO)) = LID}.
As validações locais são:
\begin{lstlisting}
eq valido?(livro(LID, OID, T, AUT, ANO)) =
  (T =/= "") and (AUT =/= "") and (ANO > 0) .
eq valido?(usuario(UID, NOM, EMAIL)) =
  (NOM =/= "") and (EMAIL =/= "") .
\end{lstlisting}

\subsection{Acervo}

As equações abaixo representam inserção guardada, remoção de todas as
ocorrências de um ID e contagem recursiva, incluindo o caso vazio:
\begin{lstlisting}
eq adicionarLivro(L, A) =
  if not(pertenceAcervo?(id(L), A)) and valido?(L)
  then L A else A fi .
eq removerDoAcervo(LID, vazioAcervo) = vazioAcervo .
eq removerDoAcervo(LID, L A) =
  if LID == id(L) then removerDoAcervo(LID, A)
  else L removerDoAcervo(LID, A) fi .
eq qtdAcervo(vazioAcervo) = 0 .
eq qtdAcervo(L A) = 1 + qtdAcervo(A) .
\end{lstlisting}

A busca filtra o ID e distingue ausência, resultado único e duplicidade:
\begin{lstlisting}
eq buscarLivro(LID, A) = resultadoBusca(filtrarId(LID, A)) .
eq filtrarId(LID, vazioAcervo) = vazioAcervo .
eq filtrarId(LID, L A) =
  if LID == id(L) then L filtrarId(LID, A)
  else filtrarId(LID, A) fi .
eq resultadoBusca(vazioAcervo) = livroInexistente .
eq resultadoBusca(L) = L .
eq resultadoBusca(L L' A) = livroDuplicado .
\end{lstlisting}

\texttt{buscarTitulo} e \texttt{tituloDeAutor} também filtram o acervo, por
igualdade de título ou autor; devolvem \texttt{vazioAcervo} quando não há
correspondência. \texttt{qtdExemplaresObra} conta os livros com o
\texttt{ObraId} informado. A coleção de usuários segue a mesma recursão de
inserção, remoção e contagem, com \texttt{vazioUsuarios} como caso base e
\texttt{pertence?} para verificar o ID; a inserção também exige
\texttt{valido?(\allowbreak{}U)}. A remoção de um ID ausente preserva a coleção.

\subsection{Empréstimos e devoluções}

Na coleção de empréstimos, a inserção evita repetir o mesmo par livro--usuário;
a devolução remove os registros do exemplar. As restrições relativas ao acervo
e ao cadastro são aplicadas pelas transições globais.
\begin{lstlisting}
eq emprestar(LID, UID, ES) =
  if temLivro?(UID, LID, ES) then ES
  else emprestar(LID, UID) ES fi .
eq devolver(LID, semEmprestimo) = semEmprestimo .
eq devolver(LID, emprestar(LID2, UID) ES) =
  if LID == LID2 then devolver(LID, ES)
  else emprestar(LID2, UID) devolver(LID, ES) fi .
eq qtdEmprestimos(semEmprestimo) = 0 .
eq qtdEmprestimos(emprestar(LID, UID) ES) = 1 + qtdEmprestimos(ES) .
\end{lstlisting}

\texttt{emprestado?} procura o ID do livro; \texttt{temLivro?} procura um par
usuário--livro, ou qualquer empréstimo do usuário na versão de dois argumentos.
Esses predicados retornam \texttt{false} sobre a coleção vazia.

\subsection{Transições globais}

As seis transições recebem um estado e produzem outro. Os registros delegam às
inserções guardadas das coleções; as remoções exigem livro disponível ou usuário
sem empréstimos. Emprestar exige disponibilidade e cadastro, enquanto devolver
exige que o livro esteja com o usuário informado.
\begin{lstlisting}
eq registrarLivro(L, biblioteca(A, US, ES)) =
  biblioteca(adicionarLivro(L, A), US, ES) .
eq registrarUsuario(U, biblioteca(A, US, ES)) =
  biblioteca(A, adicionarUsuario(U, US), ES) .
eq removerLivro(LID, biblioteca(A, US, ES)) =
  if disponivel?(LID, biblioteca(A, US, ES))
  then biblioteca(removerDoAcervo(LID, A), US, ES)
  else biblioteca(A, US, ES) fi .
eq removerUsuario(UID, biblioteca(A, US, ES)) =
  if not(temLivro?(UID, ES))
  then biblioteca(A, removerUsuario(UID, US), ES)
  else biblioteca(A, US, ES) fi .
eq emprestarLivro(LID, UID, biblioteca(A, US, ES)) =
  if disponivel?(LID, biblioteca(A, US, ES)) and
     usuarioRegistrado?(UID, biblioteca(A, US, ES))
  then biblioteca(A, US, emprestar(LID, UID, ES))
  else biblioteca(A, US, ES) fi .
eq devolverLivro(LID, UID, biblioteca(A, US, ES)) =
  if temLivro?(UID, LID, biblioteca(A, US, ES))
  then biblioteca(A, US, devolver(LID, ES))
  else biblioteca(A, US, ES) fi .
\end{lstlisting}

Os observadores globais consultam as coleções pelas projeções. Por exemplo,
\texttt{usuarioRegistrado?(\allowbreak{}UID,\allowbreak{} B)} reduz a
\texttt{pertence?(\allowbreak{}UID,\allowbreak{} usuarios(\allowbreak{}B))}, e \texttt{buscarTitulo(\allowbreak{}T,\allowbreak{} B)} reduz a
\texttt{buscarTitulo(\allowbreak{}T,\allowbreak{} acervo(\allowbreak{}B))}. A disponibilidade e sua contagem são:
\begin{lstlisting}
eq disponivel?(LID, B) =
  pertenceAcervo?(LID, B) and not emprestado?(LID, B) .
eq qtdDisponiveis(B) = contarDisponiveis(acervo(B), emprestimos(B)) .
eq contarDisponiveis(vazioAcervo, ES) = 0 .
eq contarDisponiveis(L A, ES) =
  if emprestado?(id(L), ES) then contarDisponiveis(A, ES)
  else 1 + contarDisponiveis(A, ES) fi .
\end{lstlisting}

\subsection{Invariantes globais}

\texttt{idsUnicos?} verifica ausência de IDs repetidos no acervo e nos usuários;
\texttt{todosValidos?} verifica os atributos de cada entidade.
\texttt{livrosUnicos?} impede que um exemplar apareça em mais de um registro
de empréstimo. Esses predicados retornam \texttt{true} nas coleções vazias e
percorrem as coleções não vazias. A consistência das referências e a validade
global são definidas por:
\begin{lstlisting}
eq emprestimosValidos?(biblioteca(A, US, semEmprestimo)) = true .
eq emprestimosValidos?(biblioteca(A, US, emprestar(LID, UID) ES)) =
  pertenceAcervo?(LID, A) and pertence?(UID, US) and
  emprestimosValidos?(biblioteca(A, US, ES)) .
eq valido?(biblioteca(A, US, ES)) =
  idsUnicos?(A) and todosValidos?(A) and
  idsUnicos?(US) and todosValidos?(US) and
  livrosUnicos?(ES) and emprestimosValidos?(biblioteca(A, US, ES)) .
\end{lstlisting}

Em estados válidos, cada exemplar pertence a exatamente uma das categorias
disponível ou emprestado, e cada empréstimo corresponde a um exemplar do
acervo. Portanto, \texttt{qtdDisponiveis(\allowbreak{}B) + qtdEmprestimos(\allowbreak{}B) = qtdAcervo(\allowbreak{}B)}.
Os predicados permitem detectar estados inconsistentes construídos diretamente
com \texttt{biblioteca}; as transições preservam a validade quando partem de
um estado válido.

\section{Tratamento de estados inválidos}

\subsection{Busca por livro inexistente}

\subsection{Busca por livro com ID duplicado}

\subsection{Remoção de entidade inexistente}

\subsection{Adição de elemento inválido ou com ID inexistente}

\subsection{Devolução de livro não emprestado}

\subsection{Estados inválidos em nível global}

\section{Termos de teste}

Conforme os requisitos do trabalho, foram desenvolvidos os termos de teste
agregados na Tabela~\ref{tab:termos}, cobrindo as operações básicas e compostas,
casos inválidos e aplicação das regras de negócio.

Os exemplos usam o contexto abaixo, aberto após carregar
\texttt{biblioteca.mod}. \texttt{bLivre} contém um livro disponível e dois
usuários; \texttt{bEmprestada} contém o mesmo livro emprestado a Ana.
\texttt{l2} e \texttt{ux} não estão cadastrados. Cada linha da tabela é
avaliada independentemente, sem modificar os estados usados pelas demais.
\begin{lstlisting}
open BIBLIOTECA .
ops l1 l2 : -> LivroId .
ops u1 u2 ux : -> UsuarioId .
op o1 : -> ObraId .
op livro1 : -> Livro .
ops ana bia : -> Usuario .
ops bLivre bEmprestada : -> Biblioteca .
eq livro1 = livro(l1, o1, "T1", "A1", 2020) .
eq ana = usuario(u1, "Ana", "ana@x") .
eq bia = usuario(u2, "Bia", "bia@x") .
eq bLivre = biblioteca(livro1, ana bia, semEmprestimo) .
eq bEmprestada = emprestarLivro(l1, u1, bLivre) .
\end{lstlisting}

\begin{table}[htbp]
    \centering
    \small
    \renewcommand{\arraystretch}{1.2}
    \caption{Termos de teste e resultados esperados.}
    \label{tab:termos}
    \begin{tabularx}{\linewidth}{@{} c Y l @{}}
        \toprule
        \# & Termo de teste                                                                             & Resultado esperado                           \\
        \midrule
        1  & \texttt{qtdAcervo(\allowbreak{}adicionarLivro(\allowbreak{}livro1,\allowbreak{} vazioAcervo))}                                   & \texttt{1}                                   \\
        2  & \texttt{pertenceAcervo?(\allowbreak{}l1,\allowbreak{} adicionarLivro(\allowbreak{}livro1,\allowbreak{} vazioAcervo))}                         & \texttt{true}                                \\
        3  & \texttt{buscarLivro(\allowbreak{}l2,\allowbreak{} adicionarLivro(\allowbreak{}livro1,\allowbreak{} vazioAcervo))}                             & \texttt{livroInexistente}                    \\
        4  & \texttt{qtdUsuarios(\allowbreak{}adicionarUsuario(\allowbreak{}ana,\allowbreak{} vazioUsuarios))}                               & \texttt{1}                                   \\
        5  & \texttt{emprestado?(\allowbreak{}l1,\allowbreak{} emprestar(\allowbreak{}l1,\allowbreak{} u1))}                                                & \texttt{true}                                \\
        6  & \texttt{temLivro?(\allowbreak{}u1,\allowbreak{} l1,\allowbreak{} emprestar(\allowbreak{}l1,\allowbreak{} u1))}                                              & \texttt{true}                                \\
        7  & \texttt{disponivel?(\allowbreak{}l1,\allowbreak{} emprestarLivro(\allowbreak{}l1,\allowbreak{} u1,\allowbreak{} bLivre))}                                  & \texttt{false}                               \\
        8  & \texttt{qtdEmprestimos(\allowbreak{}emprestarLivro(\allowbreak{}l1,\allowbreak{} u1,\allowbreak{} bLivre))}                                   & \texttt{1}                                   \\
        9  & \texttt{emprestarLivro(\allowbreak{}l1,\allowbreak{} ux,\allowbreak{} bLivre)}                                                   & \texttt{bLivre}                              \\
        10 & \texttt{removerLivro(\allowbreak{}l1,\allowbreak{} emprestarLivro(\allowbreak{}l1,\allowbreak{} u1,\allowbreak{} bLivre))}                                 & \texttt{bEmprestada}                         \\
        11 & \texttt{devolverLivro(\allowbreak{}l1,\allowbreak{} u2,\allowbreak{} bEmprestada)}                                               & \texttt{bEmprestada}                         \\
        12 & \texttt{qtdDisponiveis(\allowbreak{}bEmprestada) + qtdEmprestimos(\allowbreak{}bEmprestada)}                         & \texttt{1}                                   \\
        \bottomrule
    \end{tabularx}
\end{table}

Para reproduzir cada verificação, usa-se \texttt{red termo == resultado .}
no contexto acima e, ao final, \texttt{close}. No teste 12, o total 1 também
corresponde a \texttt{qtdAcervo(\allowbreak{}bEmprestada)}. Nos testes 9 e 11, as guardas
recusam, respectivamente, o usuário não cadastrado e a devolução por outro
usuário. No teste 10, o estado preservado é o posterior ao empréstimo.

As sequências abaixo resumem a reescrita, com as condições já avaliadas.
No primeiro exemplo, a inserção aceita o livro válido e a contagem chega a 1:
\begin{lstlisting}
qtdAcervo(adicionarLivro(livro1, vazioAcervo))
  => qtdAcervo(livro1 vazioAcervo)
  => 1 + qtdAcervo(vazioAcervo)
  => 1 + 0 => 1
\end{lstlisting}
No segundo, o livro emprestado não pode ser removido:
\begin{lstlisting}
removerLivro(l1, emprestarLivro(l1, u1, bLivre))
  => removerLivro(l1, biblioteca(livro1, ana bia, emprestar(l1, u1)))
  => biblioteca(livro1, ana bia, emprestar(l1, u1))
\end{lstlisting}
A condição \texttt{disponivel?} é falsa nesse estado; o resultado é a forma
expandida de \texttt{bEmprestada}. O símbolo \texttt{=>} apenas indica passos
de reescrita nesta exposição, não um comando de CafeOBJ.

A suíte \texttt{spec/\allowbreak{}testes.cafe}, executada por
\texttt{python3 scripts/\allowbreak{}testar.py --homebrew} no ambiente utilizado, apresentou
\texttt{165/\allowbreak{}165} asserções verdadeiras, com 165 reduções. Dessas, 64 pertencem
ao arquivo \texttt{teste-\allowbreak{}transicoes-\allowbreak{}biblioteca.cafe}, que cobre transições,
consultas e estados inconsistentes. Esse resultado verifica os casos testados.

O cenário separado \texttt{spec/\allowbreak{}demo.cafe} começa com cinco exemplares e quatro
usuários. Após empréstimos, uma tentativa de empréstimo repetido, uma devolução
e o cadastro de outro exemplar, \texttt{biblioteca1} possui seis livros, três
empréstimos e três disponíveis. A demonstração também consulta o catálogo,
exercita operações recusadas e verifica \texttt{valido?}; ela não integra as
165 asserções da suíte.

\section{Análise da especificação}

\subsection{Justificativa dos \textit{sorts} e operadores}

Os \textit{sorts} \texttt{LivroId,\allowbreak{} UsuarioId e ObraId} foram separados para
garantir unicidade e evitar confusão entre identificadores de naturezas
distintas. \texttt{Livro} e \texttt{Usuario} são entidades com construtores e
projeções, permitindo validação (\texttt{valido?}) e consultas. As coleções
foram modeladas como monóides comutativos, já que a ordem dos elementos não é
relevante para as consultas, e a composição é naturalmente associativa e
comutativa. O \textit{sort} \texttt{Biblioteca} consolida o estado global e
permite transições atômicas e guardadas.

\subsection{Representatividade do modelo}

A especificação representa as principais características do sistema: acervo,
usuários, empréstimos, disponibilidade, pertinência, contagens e regras de
negócio. As transições globais refletem o ciclo de vida de um livro (registro,
empréstimo, devolução, remoção) e de um usuário (registro, remoção
condicionada). As invariantes \texttt{valido?} e \texttt{emprestimosValidos?}
capturam a consistência global.

\subsection{Características não representadas}

Algumas características do mundo real não foram modeladas, por não serem
exigidas pelo escopo mínimo:

\begin{enumerate}
    \item datas de empréstimo e devolução, prazos e multas;
    \item limites de quantidade de empréstimos por usuário;
    \item reservas e filas de espera;
    \item histórico de empréstimos encerrados;
    \item informações bibliográficas mais específicas (editora, edição etc).
\end{enumerate}

Essas extensões poderiam ser incorporadas com novos \textit{sorts} e equações,
sem alterar a estrutura das coleções.

\subsection{Propriedades matemáticas expressas}

A especificação expressa:

\begin{itemize}
    \item Monóides comutativos para acervo, usuários e empréstimos, com
          identidade e associatividade;

    \item Idempotência de filtros e remoções quando o elemento não existe;
    \item Conservação de cardinalidade:
          \texttt{qtdDisponiveis(\allowbreak{}B) + qtdEmprestimos(\allowbreak{}B) = qtdAcervo(\allowbreak{}B);}

    \item Unicidade de identificadores via \texttt{idsUnicos?};

    \item Validade global via \texttt{valido?}.
\end{itemize}

\subsection{Representação e mitigação de operações inválidas}

As operações inválidas são tratadas de maneiras distintas nas diferentes partes
do sistema:

\begin{itemize}
    \item Consultas ao acervo: No módulo \texttt{acervo.mod}, é utilizado o
          \textit{supersort} \texttt{LivroOuErro} que inclui dois construtores
          de erro: \texttt{livroInexistente}, retornado pela busca a um livro
          cujo ID não existe no acervo; e \texttt{livroDuplicado}), se o ID
          consultado ocorre múltiplas vezes; Essa decisão permite distinguir a
          ausência, presença ou quebra da unicidade de um livro a partir de um
          tipo bem definido.

    \item Guardas nos modificadores de coleção: Os modificadores são guardados
          por equações que, mediante falha na pré-condição, retornam a coleção
          original. Desta forma, todas as operações são totais a nível de
          coleção. Este comportamento foi implementado nas operações.
          \begin{itemize}
              \item Adição de livros e usuários: Insere apenas se o objeto -
                    livro ou usuário - é válido e se seu ID já não pertence à
                    coleção; se falhar nestas condições, devolve a coleção.

              \item Remoção de livros e usuários: Se o ID do objeto não existe,
                    retorna o estado atual da coleção. Essa abordagem torna a
                    operação idempotente.

              \item Devolução de livros: preserva a coleção de empréstimos
                    empréstimo não está presente.
          \end{itemize}

    \item Guardas nas transições globais: No módulo de \texttt{Biblioteca}, as
          transições globais são verificadas por meio de observadores,
          \texttt{disponivel?,\allowbreak{} usuarioRegistrado? e temLivro?}; se não são
          satisfeitos, a bibioteca original é retornada. Tal como nos
          modificadores de coleção, esta abordagem faz o sistema total e seguro,
          pois a partir de um estado válido, nenhuma transação levaria a um
          inválido.
\end{itemize}

A suíte de testes verifica que, após qualquer tentativa inválida, o estado
permanece inalterado e os invariantes continuam válidos.
\end{document}
